% EM structure and teaching revision 20260928
\subsection{用连续性检验位移电流项}
对安培--麦克斯韦方程取散度，利用旋度的散度为零，可得
\[
 0=\nabla\cdot\boldsymbol j+\frac{\partial}{\partial t}(\nabla\cdot\boldsymbol D)
 =\nabla\cdot\boldsymbol j+\frac{\partial\rho_f}{\partial t}.
\]
这正是自由电荷的连续性方程。位移电流不是为了给真空人为安排真实载流子，而是使变化电场、磁场与电荷守恒能够在同一组场方程中成立。
对正在充电的理想平行板电容器，取同一条闭合环路，可以选穿过导线的曲面，也可以选绕开导线、穿过极板间隙的曲面。前者包围传导电流，后者包围电位移通量的变化率，两者应给出相同环量。
若只取半径 $r$ 小于极板半径 $a$ 的同心圆回路，并忽略边缘场，极板间 $D$ 均匀，故包围的位移电流为总充电电流的 $r^2/a^2$。磁场因而随 $r$ 线性增加，而不是处处按整根导线的 $1/r$ 规律变化。曲面选择与包围比例，是这一类题目最容易混淆的两个步骤。

